\chapter{Introduction}
\label{chap:introduction}

\section{Motivation}

Between September 2019 and July 2022 the ECB held its deposit facility rate at minus 0.5 percent. Fourteen months later the rate stood at 4 percent, and by June 2025 it had been cut back to 2 percent. A path like that raises the question this thesis asks: through which channels does ECB policy reach the economy, and how long does its trace last? The question is old, but the post-COVID period sharpened it. Monetary transmission is not only about the average response of inflation or output. It is also about the channels through which financial conditions, bank intermediation, housing credit, and compensation-linked variables adjust over time. The Euro Area entered the pandemic after a long period of low rates, unconventional policy, and balance-sheet expansion. It then experienced extraordinary accommodation, followed by a rapid tightening cycle as inflation accelerated. In that environment, the same ECB announcement can move market prices inside the high-frequency window, shift lending margins over subsequent months, affect house-purchase credit growth, and leave a weaker or less consistent trace in compensation-linked proxies.

The stakes of that question are visible in the data this thesis uses. Between the first quarter of 2015 and the end of 2021, German house prices deflated by consumer prices rose by about 49 percent, while real compensation per employee in the euro area rose by about 4 percent over the whole decade to 2025. House-purchase lending in the euro area grew by about 3.3 percent a year through the low-rate years, slowed to about 2.2 percent a year once policy tightened, and dipped just below zero in April 2024. The gap between house prices and pay is the affordability backdrop against which the question is asked. This thesis does not measure who gained or lost access to housing. It measures the step before that: whether, and for how long, ECB policy news reaches housing credit and lending conditions at all, and whether it reaches them more clearly than it reaches pay.

This thesis studies that heterogeneous transmission. Its central argument is not that monetary policy mainly inflated assets or that broad real-economy transmission was absent. Those claims exceed what aggregate local projections can identify. The thesis instead asks whether ECB easing shocks, read against the post-COVID policy cycle, transmitted most clearly and persistently through housing-credit and selected financial-intermediation variables, while evidence for compensation-linked real-economy transmission was weaker and less consistent.

The empirical contribution is therefore mechanism-focused. The required interpretation chain is:

\begin{figure}[htbp]
  \centering
  \caption{Reduced-form transmission chain guiding the empirical interpretation. The diagram is a mechanism guide, not a structurally identified mediation model. Equity and asset-market responses are read alongside the chain under information-effect caution, not as a strict downstream consequence of housing-credit expansion.}
  \label{fig:conceptual-intermediation-chain}
  \thesismechanismchain
    {ECB policy-news shock}
    {lending conditions and intermediation variables}
    {house-purchase credit growth}
    {selected financial and asset-market responses}
  \figuresource{own illustration.}
\end{figure}

\noindent The framework is a high-frequency identified, dynamic impulse-response design. Its focus is interpretive. Housing-credit transmission and intermediation variables sit at the center of the thesis, while equity prices and broad financial indices are supporting evidence that must be read with information-effect caution.

\section{Research Question}

The thesis is organized around the following research question:

\begin{quote}
Did post-COVID ECB monetary easing transmit more strongly through housing-credit and financial-intermediation channels than through compensation-linked real-economy variables?
\end{quote}

The question is comparative, dynamic, and reduced-form. It concerns relative response strength and persistence across observable channels. Stronger transmission here means greater persistence, an accumulated response that keeps its sign and excludes zero through the longer horizons, together with sign consistency across horizons and samples. It does not mean significance at a single medium horizon alone. The question does not ask whether monetary policy caused household-level welfare changes, generated inequality, or shifted resources across population groups. Those questions require household, regional, or bank-level microdata that are outside the scope of the empirical design.

\clearpage
\section{Hypotheses}

Three hypotheses discipline the empirical interpretation.

\begin{hypothesis}[Housing-credit persistence]
\label{hyp:housing-credit-persistence}
Expansionary ECB policy-news shocks generated clear and persistent responses in housing-credit variables, especially house-purchase lending growth.
\end{hypothesis}

\begin{hypothesis}[Compensation-linked attenuation]
\label{hyp:compensation-linked-attenuation}
Compensation-linked real-economy variables responded less consistently and less persistently than the housing-credit block.
\end{hypothesis}

\begin{hypothesis}[Differential financial-intermediation transmission]
\label{hyp:differential-financial-intermediation}
Expansionary monetary-policy shocks generated stronger and more persistent responses in housing-credit and financial-intermediation variables than in productive-credit and compensation-linked variables, indicating heterogeneous monetary transmission across sectors.
\end{hypothesis}

These hypotheses keep the thesis close to the evidence. \Cref{hyp:housing-credit-persistence} gives house-purchase lending growth its central role. \Cref{hyp:compensation-linked-attenuation} frames wages and labor-income proxies as a comparative block rather than as a welfare test. \Cref{hyp:differential-financial-intermediation} replaces overbroad distributional language with a testable macro-financial transmission claim. \Cref{sec:hypothesis-derivation} derives all three from the credit channel of monetary transmission and states, for each, the pattern in the response paths that would reject it.

\section{Empirical Object}

The thesis estimates monthly local-projection responses to high-frequency ECB target-factor surprises, aggregated to the month and signed so that positive values denote easing. Responses are estimated at horizons of 0, 1, 3, 6, 12, and 24 months and normalized to a one-standard-deviation surprise. Each response is estimated on all months available for it between 2005 and 2025, for statistical power. The label post-COVID denotes the motivating policy context that renews the salience of these channels, not a restriction of the estimation window.

The empirical object is a horizon-specific response profile to measured ECB policy news. The comparison spans house-purchase lending growth, pure-new mortgage flows, mortgage and NFC lending spreads, broad credit stocks, ECB liquidity-stock context, equity-market responses, labor expectations, and negotiated wage-pressure proxies. Quarterly house prices, compensation per employee, and Bank Lending Survey indicators are not interpolated into monthly responses. They remain context or robustness material when their observed frequency allows.

This object supports claims about differential transmission strength, response persistence, timing, and robustness. It does not identify household-level incidence, welfare effects, inequality, redistribution, exact structural mediation, or bank-specific treatment effects. That boundary is central to the thesis because the final inference is strongest when the estimand remains tied to observable repository outputs.

\section{Contribution}

The thesis makes five contributions.

\begin{enumerate}
  \item It brings high-frequency ECB policy-news surprises into a monthly local-projection framework and traces their effects across housing credit, lending conditions, credit aggregates, asset prices, and compensation-linked variables over two years.

  \item It makes house-purchase lending growth the central empirical variable and shows that it carries the clearest and most persistent response to policy news, a housing-credit result rather than a broad asset-market claim.

  \item It measures financial intermediation directly, through lending margins over the policy rate, lending rates, credit aggregates, balance-sheet context, and timing evidence, and shows that the margin response is the most stable in the whole set.

  \item It screens every announcement for information effects and reads the negative real DAX response through that screen, as mixed financial evidence rather than as a contradiction.

  \item It states the claim boundary directly and reports its own uncertainty in full, including where the accumulated intervals depend on how dependence across horizons is handled, so that every claim can be traced to a reproducible output.
\end{enumerate}
